% =============================================================================
%                    WEEK 3: STATISTICAL LANGUAGE MODELS
% =============================================================================
\section{Week 3: Statistical Language Models}

% -----------------------------------------------------------------------------
%                              3.1 TOPIC SUMMARY
% -----------------------------------------------------------------------------
\subsection{Topic Summary}

\begin{keybox}[Learning Objectives]
By the end of this week, you will be able to:
\begin{enumerate}
  \item Define a \textbf{language model} and use the \textbf{chain rule} with the \textbf{Markov assumption} to build N-gram approximations.
  \item Estimate N-gram probabilities with \textbf{maximum likelihood estimation (MLE)} from corpus counts.
  \item Handle sparsity via \textbf{smoothing}, \textbf{interpolation}, and \textbf{backoff}.
  \item Evaluate language models using \textbf{training/dev/test sets} and \textbf{perplexity}.
  \item Explain strengths and limitations of N-gram models, including their inability to capture \textbf{long-distance dependencies}.
\end{enumerate}
\end{keybox}

\separator

\subsubsection*{The Big Picture}

This week answers: \textit{``How can we predict the next word in a sequence?''}

\textbf{N-gram language models} assign probabilities to word sequences by:
\begin{itemize}
  \item Decomposing $P(w_1, \ldots, w_n)$ using the \textbf{chain rule}.
  \item Approximating long histories with the \textbf{Markov assumption} (only last $N-1$ words matter).
  \item Estimating probabilities from \textbf{corpus counts} (MLE).
\end{itemize}

\textbf{Why study N-grams?} They introduce foundational concepts (perplexity, smoothing, train/test splits) that carry over to modern LLMs.

\separator

\subsubsection*{What Examiners Like to Ask}

\begin{itemize}
  \item Define a language model; write the chain rule decomposition.
  \item State the Markov assumption for bigrams/trigrams.
  \item Compute bigram probabilities from counts (MLE).
  \item Explain why zeros occur and how \textbf{add-one smoothing} addresses them.
  \item Define \textbf{perplexity}; explain why lower is better.
  \item Compare \textbf{interpolation} vs.\ \textbf{backoff}.
  \item Explain why N-grams fail on long-distance dependencies.
\end{itemize}

\separator

% -----------------------------------------------------------------------------
%                 3.2 OVERVIEW + CORE CONCEPTS + EXTENDED THEORY
% -----------------------------------------------------------------------------
\subsection{Overview + Core Concepts + Extended Theory}

\subsubsection*{Roadmap}
\begin{enumerate}
  \item What is a language model?
  \item Chain rule and Markov assumption
  \item N-gram models: unigram, bigram, trigram
  \item Maximum likelihood estimation (MLE)
  \item Practical issues: log probabilities, vocabulary
  \item Evaluation: train/dev/test sets and perplexity
  \item Sampling from N-gram models
  \item Sparsity and zeros
  \item Smoothing: add-one (Laplace)
  \item Interpolation and backoff
  \item Limitations of N-gram models
\end{enumerate}

\bigseparator

% --- What is a Language Model? ---
\subsubsection{What is a Language Model?}

\begin{defbox}[Language Model (LM)]
A \textbf{language model} assigns probabilities to sequences of words.

\textbf{Two equivalent views:}
\begin{enumerate}
  \item Probability of a whole sentence: $P(w_1, w_2, \ldots, w_n)$
  \item Probability of the next word given history: $P(w_n | w_1, \ldots, w_{n-1})$
\end{enumerate}
\end{defbox}

\textbf{Applications:}
\begin{itemize}
  \item \textbf{Spell/grammar checking:} ``Their are two midterms'' $\to$ ``There are two midterms''
  \item \textbf{Speech recognition:} ``I will be back soonish'' vs.\ ``I will be bassoon dish''
  \item \textbf{Text generation:} LLMs predict words left-to-right (autoregressive).
\end{itemize}

\separator

% --- Chain Rule ---
\subsubsection{The Chain Rule of Probability}

\begin{defbox}[Chain Rule]
For any sequence of random variables:
\[
P(x_1, x_2, \ldots, x_n) = \prod_{k=1}^{n} P(x_k | x_1, \ldots, x_{k-1})
\]

Applied to words:
\[
P(w_1, \ldots, w_n) = P(w_1) \cdot P(w_2|w_1) \cdot P(w_3|w_1, w_2) \cdots P(w_n|w_1, \ldots, w_{n-1})
\]
\end{defbox}

\begin{exbox}[Example: Chain Rule]
$P(\text{``The water of Walden Pond''})$
\[
= P(\text{The}) \times P(\text{water}|\text{The}) \times P(\text{of}|\text{The water}) \times \cdots
\]
\end{exbox}

\textbf{Problem:} We can't estimate $P(w_n | w_1, \ldots, w_{n-1})$ directly --- there are too many possible histories!

\separator

% --- Markov Assumption ---
\subsubsection{The Markov Assumption}

\begin{defbox}[Markov Assumption]
Approximate the full history by only the last $N-1$ words:
\[
P(w_n | w_1, \ldots, w_{n-1}) \approx P(w_n | w_{n-N+1}, \ldots, w_{n-1})
\]

\textbf{Intuition:} Recent context is most predictive; distant history matters less.
\end{defbox}

\separator

% --- N-gram Models ---
\subsubsection{N-gram Models}

\begin{defbox}[N-gram Model]
An \textbf{N-gram model} conditions on the previous $N-1$ words.

\begin{center}
\renewcommand{\arraystretch}{1.3}
\begin{tabular}{@{} l c l @{}}
\toprule
\textbf{Model} & \textbf{N} & \textbf{Approximation} \\
\midrule
Unigram & 1 & $P(w_n)$ --- no context \\
Bigram & 2 & $P(w_n | w_{n-1})$ --- 1 word of context \\
Trigram & 3 & $P(w_n | w_{n-2}, w_{n-1})$ --- 2 words of context \\
4-gram & 4 & $P(w_n | w_{n-3}, w_{n-2}, w_{n-1})$ \\
\bottomrule
\end{tabular}
\end{center}
\end{defbox}

\begin{exbox}[Example: Bigram Approximation]
Instead of:
\[
P(\text{blue} | \text{The water of Walden Pond is so beautifully})
\]
Approximate with:
\[
P(\text{blue} | \text{beautifully})
\]
\end{exbox}

\begin{exbox}[Example: Shakespeare Generation by N-gram Order]
\textbf{Unigram:} ``To him swallowed confess hear both. Which. Of save on trail for are ay device\ldots''

\textbf{Bigram:} ``Why dost stand forth thy canopy, forsooth; he is this palpable hit the King Henry.''

\textbf{Trigram:} ``Fly, and will rid me these news of price. Therefore the sadness of parting, as they say, 'tis done.''

\textbf{4-gram:} ``King Henry. What! I will go seek the traitor Gloucester.''

\textbf{Observation:} Higher N $\Rightarrow$ more coherent text (but more memorisation of training data).
\end{exbox}

\separator

% --- MLE ---
\subsubsection{Maximum Likelihood Estimation (MLE)}

\begin{defbox}[MLE for Bigrams]
\[
P(w_n | w_{n-1}) = \frac{C(w_{n-1}, w_n)}{C(w_{n-1})}
\]

where $C(\cdot)$ denotes corpus count.

\textbf{General N-gram:}
\[
P(w_n | w_{n-N+1}^{n-1}) = \frac{C(w_{n-N+1}^{n})}{C(w_{n-N+1}^{n-1})}
\]
\end{defbox}

\begin{exbox}[Example: Bigram Counts (Berkeley Restaurant Project)]
From 9,222 sentences:

\begin{center}
\renewcommand{\arraystretch}{1.2}
\begin{tabular}{@{} l c c c c @{}}
\toprule
 & \code{want} & \code{to} & \code{eat} & \code{chinese} \\
\midrule
\code{i} & 827 & 0 & 9 & 0 \\
\code{want} & 0 & 608 & 1 & 6 \\
\code{to} & 0 & 4 & 686 & 2 \\
\bottomrule
\end{tabular}
\end{center}

Given unigram counts: $C(\text{i}) = 2533$, $C(\text{want}) = 927$, $C(\text{to}) = 2417$

\[
P(\text{want} | \text{i}) = \frac{827}{2533} = 0.33, \quad P(\text{to} | \text{want}) = \frac{608}{927} = 0.66
\]
\end{exbox}

\begin{exbox}[Example: Sentence Probability]
$P(\text{<s> i want english food </s>})$
\[
= P(\text{i}|\text{<s>}) \times P(\text{want}|\text{i}) \times P(\text{english}|\text{want}) \times P(\text{food}|\text{english}) \times P(\text{</s>}|\text{food})
\]
\[
= 0.25 \times 0.33 \times 0.0011 \times 0.5 \times 0.68 \approx 0.000031
\]
\end{exbox}

\begin{tipbox}[What N-grams Capture]
\begin{itemize}
  \item \textbf{Grammar:} $P(\text{to}|\text{want}) = 0.66$ (want $\to$ infinitive)
  \item \textbf{World knowledge:} $P(\text{chinese}|\text{want}) > P(\text{english}|\text{want})$
  \item \textbf{Semantics:} ``dinner'' and ``lunch'' have similar distributions
\end{itemize}
\end{tipbox}

\separator

% --- Log Probabilities ---
\subsubsection{Practical Issues: Log Probabilities}

\begin{defbox}[Log Probabilities]
Multiplying many small probabilities $\Rightarrow$ \textbf{underflow}.

\textbf{Solution:} Work in log space.
\[
\log(p_1 \times p_2 \times \cdots \times p_n) = \log p_1 + \log p_2 + \cdots + \log p_n
\]

Convert back at the end: $p = \exp(\text{logprob})$
\end{defbox}

\separator

% --- Sentence Boundaries ---
\subsubsection{Sentence Boundary Tokens}

\begin{defbox}[Start and End Tokens]
Insert special tokens to model sentence boundaries:
\begin{itemize}
  \item \code{<s>} --- start of sentence
  \item \code{</s>} --- end of sentence
\end{itemize}

This allows the model to learn $P(w_1 | \text{<s>})$ and $P(\text{</s>} | w_n)$.
\end{defbox}

\bigseparator

% --- Evaluation ---
\subsubsection{Evaluation: Train / Dev / Test Sets}

\begin{defbox}[Data Splits]
\begin{center}
\renewcommand{\arraystretch}{1.3}
\begin{tabular}{@{} l p{9cm} @{}}
\toprule
\textbf{Set} & \textbf{Purpose} \\
\midrule
\textbf{Training} & Estimate model parameters (N-gram counts) \\
\textbf{Development (Dev)} & Tune hyperparameters (e.g., interpolation $\lambda$s); can evaluate repeatedly \\
\textbf{Test} & Final evaluation \textit{only once} (or very few times) \\
\bottomrule
\end{tabular}
\end{center}
\end{defbox}

\begin{tipbox}[Training on the Test Set = Bad Science!]
If test sentences influence training:
\begin{itemize}
  \item Model assigns artificially high probability to test data.
  \item Evaluation metric is falsely optimistic.
  \item Results don't generalise.
\end{itemize}
\end{tipbox}

\separator

% --- Perplexity ---
\subsubsection{Perplexity}

\begin{defbox}[Perplexity (PP)]
\[
\text{PP}(W) = P(w_1, w_2, \ldots, w_N)^{-1/N} = \sqrt[N]{\frac{1}{P(w_1, \ldots, w_N)}}
\]

\textbf{Interpretation:}
\begin{itemize}
  \item Inverse probability, normalised by sequence length.
  \item \textbf{Weighted average branching factor:} how many equally likely words could follow on average.
  \item \textbf{Lower perplexity = better model.}
\end{itemize}

\textbf{Range:} $[1, \infty)$ (probability in $[0, 1]$).
\end{defbox}

\begin{exbox}[Example: Perplexity Calculation]
Deterministic language $L = \{\text{red}, \text{blue}, \text{green}\}$, each word equally likely ($p = 1/3$).

Test set $T = $ ``red red red red blue'' (5 words).

\[
P(T) = (1/3)^5, \quad \text{PP}(T) = ((1/3)^5)^{-1/5} = (1/3)^{-1} = 3
\]

Now suppose training made \code{red} very likely: $P(\text{red}) = 0.8$, $P(\text{blue}) = 0.1$.

\[
P(T) = 0.8^4 \times 0.1 = 0.04096, \quad \text{PP}(T) = 0.04096^{-1/5} \approx 1.89
\]

\textbf{Lower perplexity} $\Rightarrow$ model better predicts the test set.
\end{exbox}

\begin{exbox}[Example: Perplexity by N-gram Order (WSJ)]
Training: 38M words; Test: 1.5M words.

\begin{center}
\renewcommand{\arraystretch}{1.2}
\begin{tabular}{@{} l c c c @{}}
\toprule
 & Unigram & Bigram & Trigram \\
\midrule
Perplexity & 962 & 170 & 109 \\
\bottomrule
\end{tabular}
\end{center}

\textbf{Observation:} More context $\Rightarrow$ lower perplexity.
\end{exbox}

\begin{tipbox}[Why Log Perplexity?]
Log perplexity turns products into sums, making computation \textbf{numerically stable}.
\end{tipbox}

\bigseparator

% --- Sampling ---
\subsubsection{Sampling from N-gram Models}

\begin{defbox}[Shannon Sampling Method]
\begin{enumerate}
  \item Sample $w_1$ from $P(w | \text{<s>})$.
  \item Sample $w_2$ from $P(w | w_1)$ (for bigram).
  \item Repeat until \code{</s>} is sampled.
  \item Concatenate the words.
\end{enumerate}

\textbf{Use:} Visualise what the model has learned; diagnose overfitting.
\end{defbox}

\separator

% --- Sparsity and Zeros ---
\subsubsection{Sparsity and Zeros}

\begin{defbox}[The Sparsity Problem]
\textbf{Most possible N-grams never appear in training.}

\textbf{Shakespeare corpus:} $N = 884{,}647$ tokens, $V = 29{,}066$ types.
\begin{itemize}
  \item Possible bigrams: $V^2 \approx 844$ million.
  \item Observed bigram types: $\sim 300{,}000$ (only 0.04\%).
\end{itemize}

\textbf{Consequence:} Unseen N-grams get $P = 0$ under MLE.
\end{defbox}

\begin{tipbox}[Why Zeros Are Fatal]
\begin{itemize}
  \item If any N-gram in test set has $P = 0$, the whole sentence has $P = 0$.
  \item Perplexity = $\infty$ (can't divide by 0).
\end{itemize}
\end{tipbox}

\begin{exbox}[Example: Zero Probability]
Training: ``ate lunch'', ``ate dinner'', ``ate a'', ``ate the''

Test: ``ate \textbf{breakfast}''

$\Rightarrow P(\text{breakfast} | \text{ate}) = 0$ under MLE.
\end{exbox}

\bigseparator

% --- Smoothing ---
\subsubsection{Smoothing: Add-One (Laplace)}

\begin{defbox}[Add-One Smoothing]
Pretend every N-gram was seen one extra time.

\textbf{MLE:}
\[
P(w_n | w_{n-1}) = \frac{C(w_{n-1}, w_n)}{C(w_{n-1})}
\]

\textbf{Add-one (Laplace):}
\[
P_{\text{Laplace}}(w_n | w_{n-1}) = \frac{C(w_{n-1}, w_n) + 1}{C(w_{n-1}) + V}
\]

where $V$ = vocabulary size.
\end{defbox}

\begin{exbox}[Example: Add-One Smoothing]
Raw counts:

\begin{center}
\renewcommand{\arraystretch}{1.2}
\begin{tabular}{@{} l c c c @{}}
\toprule
 & \code{want} & \code{to} & \code{eat} \\
\midrule
\code{i} & 827 & 0 & 9 \\
\bottomrule
\end{tabular}
\end{center}

With $C(\text{i}) = 2533$ and $V = 1446$:

\[
P_{\text{Laplace}}(\text{to} | \text{i}) = \frac{0 + 1}{2533 + 1446} = \frac{1}{3979} \approx 0.00025
\]

Previously zero $\Rightarrow$ now non-zero!
\end{exbox}

\begin{tipbox}[Problem with Add-One]
Add-one is a \textbf{blunt instrument}:
\begin{itemize}
  \item Shifts too much probability mass from frequent to rare/unseen events.
  \item Reconstituted counts show dramatic changes (e.g., 827 $\to$ 527).
  \item Generally \textbf{not used} for N-grams in practice.
\end{itemize}

\textbf{When it's used:} Naive Bayes text classification, domains with fewer zeros.
\end{tipbox}

\separator

% --- Interpolation and Backoff ---
\subsubsection{Interpolation and Backoff}

\begin{defbox}[Interpolation]
Combine estimates from different N-gram orders:
\[
\hat{P}(w_n | w_{n-2}, w_{n-1}) = \lambda_1 P(w_n | w_{n-2}, w_{n-1}) + \lambda_2 P(w_n | w_{n-1}) + \lambda_3 P(w_n)
\]

where $\lambda_1 + \lambda_2 + \lambda_3 = 1$.

\textbf{Key idea:} If trigram evidence is sparse, lean more on bigram/unigram.
\end{defbox}

\begin{defbox}[Backoff]
Use the highest-order N-gram with sufficient evidence:
\begin{enumerate}
  \item If $C(w_{n-2}, w_{n-1}, w_n) > 0$, use trigram.
  \item Else if $C(w_{n-1}, w_n) > 0$, use bigram.
  \item Else use unigram.
\end{enumerate}

\textbf{Complication:} Must discount higher-order N-grams so probabilities sum to 1.
\end{defbox}

\begin{defbox}[Stupid Backoff]
Simplified backoff without proper discounting (not a true probability distribution):
\[
S(w_i | w_{i-k+1}^{i-1}) = 
\begin{cases}
\dfrac{C(w_{i-k+1}^i)}{C(w_{i-k+1}^{i-1})} & \text{if } C(w_{i-k+1}^i) > 0 \\[8pt]
0.4 \cdot S(w_i | w_{i-k+2}^{i-1}) & \text{otherwise}
\end{cases}
\]

\textbf{Works well at web scale} (e.g., Google N-grams).
\end{defbox}

\begin{tipbox}[How to Set $\lambda$s?]
Tune on a \textbf{held-out (dev) set}:
\begin{enumerate}
  \item Fix N-gram probabilities from training data.
  \item Search for $\lambda$s that minimise perplexity on dev set.
\end{enumerate}
\end{tipbox}

\begin{tipbox}[Interpolation vs.\ Backoff]
\textbf{Interpolation} usually works better because it always uses information from all orders, not just the highest available.
\end{tipbox}

\bigseparator

% --- Limitations ---
\subsubsection{Limitations of N-gram Models}

\begin{tipbox}[Why N-grams Fail]
\begin{enumerate}
  \item \textbf{Long-distance dependencies:}
  
  ``The soups that I made from that new cookbook I bought yesterday \textbf{were} amazingly delicious.''
  
  Trigram can't see that ``soups'' (plural) requires ``were'' --- too far away.
  
  \item \textbf{No generalisation across similar meanings:}
  
  N-grams treat ``car'' and ``automobile'' as completely unrelated dimensions.
  
  \item \textbf{Sparsity grows exponentially with N.}
\end{enumerate}

\textbf{Solution:} Large language models (LLMs) with embeddings and attention.
\end{tipbox}

\textbf{But N-grams are still useful:}
\begin{itemize}
  \item Fast and lightweight.
  \item Good for constrained domains.
  \item Introduce key concepts: perplexity, smoothing, train/test, sampling.
\end{itemize}

\bigseparator

% --- Summary ---
\subsubsection{Summary: N-gram Language Models}

\begin{center}
\renewcommand{\arraystretch}{1.3}
\begin{tabular}{@{} l p{10cm} @{}}
\toprule
\textbf{Concept} & \textbf{Key Point} \\
\midrule
Language model & Assigns probability to word sequences \\
Chain rule & $P(w_1, \ldots, w_n) = \prod_k P(w_k | w_1, \ldots, w_{k-1})$ \\
Markov assumption & Approximate history with last $N-1$ words \\
MLE & $P(w_n | w_{n-1}) = C(w_{n-1}, w_n) / C(w_{n-1})$ \\
Perplexity & $\text{PP} = P(W)^{-1/N}$; lower = better \\
Add-one smoothing & Add 1 to all counts; blunt but simple \\
Interpolation & Weighted mix of unigram, bigram, trigram \\
Backoff & Fall back to lower-order N-gram if higher is zero \\
Limitation & Can't capture long-distance dependencies \\
\bottomrule
\end{tabular}
\end{center}


